% META: {"id": "TSPE-LIMFCT-COURS-07-TC", "chapitre": "TSPE-LIMITES-FONCTIONS", "type_objet": "cours", "sous_type": "td_contextualise", "titre": "TD contextualise — Concentration d'un medicament", "temps_gabarit": 7, "capacites": ["C1", "C2", "C3"], "mode_creation": "ex_nihilo", "sources_inspiration": [], "status": "generated"}

% BEGIN-VERIFY
% from sympy import *
% x, t = symbols('x t', positive=True)
% # C(t) = 5t * e^(-0.5t)
% C = 5*t*exp(-t/2)
% # C(0) = 0
% assert C.subs(t, 0) == 0
% # C'(t) = 5*e^(-t/2) * (1 - t/2)
% Cp = diff(C, t)
% assert simplify(Cp - 5*exp(-t/2)*(1 - t/2)) == 0
% # C'(t) = 0 at t = 2
% assert solve(1 - t/2, t) == [2]
% # C(2) = 10*e^(-1)
% assert C.subs(t, 2) == 10*exp(-1)
% # lim C(t) as t -> +inf = 0
% assert limit(C, t, oo) == 0
% # AH: y = 0
% END-VERIFY

\section*{Temps 7 — TD contextualise : Concentration d'un medicament}

\textbf{Contexte.} Apres injection, la concentration (en mg/L) d'un medicament dans le sang au temps $t$ (en heures, $t \geq 0$) est modelisee par :
\[
  C(t) = 5t\,\mathrm{e}^{-0{,}5t}.
\]

\bigskip

%---------------------------------------------------------------
% PARTIE A — Etude de la concentration
%---------------------------------------------------------------

\begin{exercice}{TSPE-LIMFCT-COURS-07-TC-EX1}{2}{15}

\textbf{Partie A — Limites et variations}

\begin{enumerate}
  \item Calculer $C(0)$. Interpreter.
  \item Determiner $\lim_{t \to +\infty} C(t)$. Quelle croissance comparee utilise-t-on ? Interpreter.
  \item La droite $y = 0$ est-elle asymptote horizontale a la courbe ? Justifier.
  \item Calculer $C'(t)$ et montrer que $C'(t) = 5\mathrm{e}^{-0{,}5t}\left(1 - \dfrac{t}{2}\right)$.
  \item En deduire le tableau de variations de $C$.
  \item Calculer la concentration maximale et le temps ou elle est atteinte.
\end{enumerate}

\end{exercice}

\begin{corrige}{TSPE-LIMFCT-COURS-07-TC-EX1}

\commentaireMarge{Appliquer les croissances comparees.}

\textbf{Question 1.} $C(0) = 5 \times 0 \times \mathrm{e}^0 = 0$. Au moment de l'injection, la concentration est nulle (le medicament n'a pas encore ete absorbe).

\textbf{Question 2.} $C(t) = 5t\,\mathrm{e}^{-t/2}$. On pose $X = t/2$ : $C(t) = 10X\,\mathrm{e}^{-X}$, et $\lim_{X \to +\infty} X\,\mathrm{e}^{-X} = 0$ par croissances comparees ($n = 1$). Donc $\lim_{t \to +\infty} C(t) = 0$. A long terme, le medicament est elimine.

\textbf{Question 3.} Oui. $\lim_{t \to +\infty} C(t) = 0$, donc la droite $y = 0$ (l'axe des abscisses) est asymptote horizontale a la courbe en $+\infty$.

\textbf{Question 4.} $C(t) = 5t \cdot \mathrm{e}^{-t/2}$. Par la formule du produit :
\[
  C'(t) = 5\,\mathrm{e}^{-t/2} + 5t \times \left(-\frac{1}{2}\right)\mathrm{e}^{-t/2} = 5\,\mathrm{e}^{-t/2}\left(1 - \frac{t}{2}\right).
\]

\textbf{Question 5.} $C'(t) = 0 \iff t = 2$ (car $\mathrm{e}^{-t/2} > 0$). $C'(t) > 0$ pour $t < 2$, $C'(t) < 0$ pour $t > 2$. La fonction $C$ est croissante sur $[0, 2]$, decroissante sur $[2, +\infty[$.

\textbf{Question 6.} Maximum en $t = 2$ : $C(2) = 10\,\mathrm{e}^{-1} \approx 3{,}68$ mg/L. La concentration maximale est atteinte 2 heures apres l'injection.

\end{corrige}

\bigskip

%---------------------------------------------------------------
% PARTIE B — Seuil therapeutique
%---------------------------------------------------------------

\begin{exercice}{TSPE-LIMFCT-COURS-07-TC-EX2}{2}{15}

\textbf{Partie B — Seuil therapeutique et asymptote}

Le medicament est efficace lorsque $C(t) \geq 1$ mg/L.

\begin{enumerate}
  \item Expliquer pourquoi l'equation $C(t) = 1$ admet exactement deux solutions $t_1 < t_2$.
  \item Donner une valeur approchee de $t_1$ et $t_2$ a l'aide de la calculatrice.
  \item Combien de temps le medicament est-il efficace ?
  \item Quelle est la signification graphique de l'asymptote $y = 0$ dans ce contexte ?
\end{enumerate}

\end{exercice}

\begin{corrige}{TSPE-LIMFCT-COURS-07-TC-EX2}

\textbf{Question 1.} $C$ est continue, croissante sur $[0, 2]$ de $C(0) = 0$ a $C(2) \approx 3{,}68 > 1$, puis decroissante sur $[2, +\infty[$ avec $C(t) \to 0$. Par le theoreme des valeurs intermediaires, $C(t) = 1$ admet exactement une solution $t_1 \in ]0, 2[$ et une solution $t_2 \in ]2, +\infty[$.

\textbf{Question 2.} A la calculatrice : $t_1 \approx 0{,}23$~h et $t_2 \approx 7{,}15$~h.

\textbf{Question 3.} Duree d'efficacite : $t_2 - t_1 \approx 6{,}92$~heures, soit environ 7 heures.

\textbf{Question 4.} L'asymptote $y = 0$ signifie que la concentration tend vers zero a long terme : le medicament est completement elimine par l'organisme.

\end{corrige}
